\documentclass[11pt]{article}
\usepackage[margin=0.8in]{geometry}
\usepackage{hyperref}
\pagestyle{empty}

\begin{document}

\noindent
March 12, 2026

\vspace{0.7cm}

\noindent
Dear Editor,

\vspace{0.4cm}

\noindent
Please consider our manuscript entitled \textbf{Universal spectral scaling of terrain-induced fatigue} for publication as an Article in \textit{Nature Communications}.

\vspace{0.4cm}

\noindent
Understanding how terrain geometry influences mechanical fatigue remains a long-standing challenge across vehicle dynamics, geomorphology, and reliability science. Although terrain roughness is known to drive dynamic loading in off-road vehicles, a general framework connecting measurable terrain statistics to fatigue damage has not yet been established. In this work we develop a spectral scaling framework that links terrain morphology directly to fatigue response. To our knowledge, this work provides the first quantitative framework linking terrain spectral structure directly to fatigue response through a closed-form scaling relationship.

\vspace{0.4cm}

\noindent
The framework shows that terrain-induced fatigue can be predicted using two spectral parameters describing terrain roughness: the amplitude parameter ($C_z$) and spectral slope ($\beta$). We derive the theoretical relationship ($\beta_t = 7 - 2D$), connecting terrain fractal dimension to spectral structure, and demonstrate how this relationship governs fatigue accumulation in vehicle systems. The theory is evaluated through large-scale numerical experiments comprising more than 1,500 synthetic terrain realizations and approximately 18,000 vehicle simulations. Across these simulations, terrain spectral structure explains roughly 95\% of the variability in fatigue response.

\vspace{0.4cm}

\noindent
We then validate the framework using real-world terrain and vehicle measurements, including LiDAR-derived terrain data from multiple geographic regions and vehicle vibration data collected across thousands of road segments. Despite large differences in vehicle parameters and terrain environments, the observed spectral relationships remain consistent, suggesting that terrain spectral structure provides a robust predictor of fatigue response.

\vspace{0.4cm}

\noindent
We believe these results establish a general physical connection between landscape structure and mechanical wear. Such a framework may enable terrain-aware predictive maintenance, route planning, and improved vehicle design in applications ranging from autonomous ground vehicles and heavy equipment to defense mobility systems.

\vspace{0.4cm}

\noindent
All simulation code and analysis scripts used in this study are available to editors and reviewers at:

\begin{center}
\url{https://github.com/GeoGizmodo/fractalFatigueGeo.git}
\end{center}

\noindent
and will be made publicly accessible upon publication.

\vspace{0.4cm}

\noindent
We respectfully suggest the following researchers as potential reviewers whose expertise spans terrain morphology, landscape scaling, and fatigue mechanics:

\begin{itemize}
\item Gregory C. Tucker, University of Colorado Boulder \\
\url{https://cires.colorado.edu/people/gregory-c-tucker}

\item Taylor Perron, Massachusetts Institute of Technology \\
\url{https://taylorperron.org}

\item Paola Passalacqua, ETH Zurich \\
\url{https://sites.google.com/site/passalacquagroup/home}

\item Emilio Martinez-Paneda, University of Oxford \\
\url{https://mechmat.web.ox.ac.uk}
\end{itemize}

\vspace{0.4cm}

\noindent
This manuscript has not been previously published and is not under consideration elsewhere. We have not had prior discussions with a Nature Communications editor regarding this work.

\vspace{0.4cm}

\noindent
\textbf{Author Information:}

\vspace{0.2cm}

\noindent
Mehdi Zare$^{1}$, Stanislava Meredith$^{1}$, Aneesh Hariharan$^{1,*}$

\vspace{0.2cm}

\noindent
$^{1}$GeoGizmodo LLC Research Division, Tacoma, Washington, USA

\vspace{0.2cm}

\noindent
$^{*}$Corresponding author: aneesh@geogizmodo.ai

\vspace{0.4cm}

\noindent
\textbf{Author Contributions:}

\vspace{0.2cm}

\noindent
All authors contributed equally to this work. M.Z., S.M., and A.H. conceived the study, developed the theoretical framework, performed simulations, analyzed data, and wrote the manuscript.

\vspace{0.4cm}

\noindent
\textbf{Acknowledgements:}

\vspace{0.2cm}

\noindent
We thank Dr. Kathryn Lindsey, Professor, Boston College and Daniel Sheeler, Software Engineer, University of Chicago for correcting earlier versions of the manuscript. This work was supported by United States Air Force contract number FA8604-25-C-B079 through the Small Business Innovation and Research program. We thank Julio Toala, DAF GeoBase Program Futures Ops Manager, for assisting with extremely useful real information regarding actual problems to be solved for the above contract.

\vspace{0.4cm}

\noindent
\textbf{Competing Interests:}

\vspace{0.2cm}

\noindent
M.Z., S.M., and A.H. are employees of GeoGizmodo LLC, which received SBIR funding for this work. The authors declare no other competing interests.

\vspace{0.4cm}

\noindent
Thank you for your consideration.

\vspace{0.4cm}

\noindent
Sincerely,

\vspace{0.3cm}

\noindent
Aneesh Hariharan \\
Corresponding Author \\
GeoGizmodo LLC \\
\href{mailto:aneesh@geogizmodo.ai}{aneesh@geogizmodo.ai}

\vspace{0.3cm}

\noindent
Mehdi Zare \\
GeoGizmodo LLC, Research Division

\vspace{0.2cm}

\noindent
Stanislava Meredith \\
GeoGizmodo LLC, Product Division

\end{document}
